% Part: counterfactuals
% Chapter: minimal-change-semantics
% Section: contraposition

\documentclass[../../../include/open-logic-section]{subfiles}

\begin{document}

\olfileid{cnt}{min}{cpo}

\olsection{প্রতিবিপরীতকরণ}

বস্তুগত ও কঠোর শর্তবচন তাদের প্রতিবিপরীত রূপের সমতুল্য।
প্রতিবাস্তব শর্তবচনের ক্ষেত্রে তা নয়। ক্রাটজারের দেওয়া
একটি উদাহরণ দেখা যাক:
\begin{quote}
  গ্যেটে যদি ১৮৩২ সালে না মারা যেতেন, তবু এখন তিনি মৃত থাকতেন।

  গ্যেটে যদি এখন মৃত না হতেন, তবে তিনি ১৮৩২ সালে মারা যেতেন।
\end{quote}
প্রথম বাক্যটি সত্য: মানুষ শত শত বছর বাঁচে না। দ্বিতীয়টি
স্পষ্টতই মিথ্যা: গ্যেটে এখন মৃত না হলে এখনও জীবিত
থাকতেন; ফলে ১৮৩২ সালে তাঁর মৃত্যু ঘটে থাকতে পারত না।

\begin{ex}\ollabel{ex:contraposition-counterex}
  গোলক অর্থতত্ত্বে প্রতিবিপরীতকরণ অবৈধ, অর্থাৎ $p
  \cif q \Entails/ \lnot q \cif \lnot p$। ধরা যাক, $p$ মানে
  ``গ্যেটে ১৮৩২ সালে মারা যাননি'' এবং $q$ মানে ``গ্যেটে এখন মৃত''।
  % BN-SRC-428: পরের পরিতৃপ্তি-বক্তব্যের সঙ্গে একই model-name।
  এটি $\mModel{M} = \tuple{W, O, V}$ মডেলে ধরা যায়, যেখানে
  $W = \{w, w_1, w_2\}$,
  % BN-SRC-427: সম্পূর্ণ O অপেক্ষক নয়, মূল জগতে O_w নির্দিষ্ট।
  $O_w = \{\{w\}, \{w, w_1\}, \{w, w_1, w_2\}\}$,
  $V(p) = \{w_1, w_2\}$ এবং $V(q) = \{w, w_1\}$। অন্য দুই
  জগতের গোলকব্যবস্থা সংজ্ঞার শর্ত মেনে যেকোনও ভাবে নেওয়া যায়। এখানে
  $w$ বাস্তব জগৎ, যেখানে গ্যেটে ১৮৩২ সালে মারা গিয়েছেন এবং
  আজও মৃত; $w_1$ নিকটবর্তী জগৎ, যেখানে তিনি ধরা যাক
  ১৮৩৩ সালে মারা গিয়েছেন এবং আজও মৃত; আর $w_2$
  দূরবর্তী জগৎ, যেখানে তিনি এখনও জীবিত।
\begin{figure}
\begin{center}
\begin{tikzpicture}[scale=.6,layerwidth=1.5]\tiny
  \spheresystem{3}
  \begin{scope}
    \clip \propositionplot{0}{.8}{50}{4.5} -- (4.5,-4.5) -- (-4.5,-4.5) -- (-4.5,4.5) -- (4.5,4.5);
    \spherelayer{3}
  \end{scope}
  \propositionintersect{0}{2}{110}{5.5}
  \proposition{0}{.8}{50}{4.5}
  \path (0,0) node[world] {$w$};
  \spherepos{0}{2}{node[world] {$w_1$}}
  \spherepos{-50}{3}{node[world] {$w_2$}}
  \spherepos{28}{3}{node {$q$}}
  \spherepos{42}{3}{node {$\lnot q$}}
  \spherepos{-55}{4}{node {$p$}}
  \spherepos{-67}{4}{node {$\lnot p$}}
\end{tikzpicture}
\caption{প্রতিবিপরীতকরণ ব্যর্থ হওয়ার প্রত্যুদাহরণ}
\ollabel{fig:contraposition}
\end{center}
\end{figure}
$p$-স্বীকারক গোলক $S = \{w, w_1\}$-এর প্রতিটি জগতে
$p \lif q$ সত্য; তাই $\mSat{M}{p \cif q}[w]$। কিন্তু
$\lnot q$-স্বীকারক গোলক $\{w, w_1, w_2\}$-এ এমন একটি জগৎ
আছে—$w_2$—যেখানে $q$ মিথ্যা ও $p$ সত্য; তাই
$\mSat/{M}{\lnot q \lif \lnot p}[w_2]$। এখানেও
এই জগৎটিই নঞর্থক অনুপক্ষ সত্য হওয়ার একমাত্র জগৎ; তাই
উক্ত একমাত্র স্বীকারক গোলক ব্যর্থ হয়ে প্রতিবিপরীত প্রতিবাস্তব
শর্তবচনটি অসত্য।
\end{ex}

\end{document}
